\documentclass[UTF8,a4paper,12pt]{ctexart}
\usepackage[margin=2.5cm]{geometry}
\usepackage{booktabs}
\usepackage{tabularx}
\usepackage{array}
\usepackage{enumitem}
\usepackage{fancyhdr}
\usepackage{float}
\usepackage{caption}
\usepackage[hidelinks]{hyperref}
\usepackage{xcolor}
\usepackage{graphicx}
\usepackage{tikz}
\usepackage{pgfplots}
\usepackage{longtable}
\pgfplotsset{compat=1.18}

\setlength{\parindent}{2em}
\setlength{\parskip}{0.35em}
\setlist{itemsep=0.25em,topsep=0.4em,parsep=0pt,partopsep=0pt}
\renewcommand{\arraystretch}{1.45}
\captionsetup[table]{position=top,skip=10pt}
\pagestyle{fancy}
\fancyhf{}
\fancyfoot[C]{\thepage}
\renewcommand{\headrulewidth}{0pt}
\newcolumntype{Y}{>{\raggedright\arraybackslash}X}
\newcommand{\code}[1]{\texttt{#1}}
\newcommand{\result}[1]{\textcolor{red!75!black}{\textbf{#1}}}
\sloppy

\title{基于 RISC-V 指令集的 CPU 设计\\后续优化说明报告}
\author{队伍编号：CICC1004252\\团队名称：和溢位}
\date{2026 年 7 月}

\begin{document}

\begin{titlepage}
  \centering
  \vspace*{2.5cm}
  {\zihao{1}\bfseries 第十届全国大学生\\[0.4em]集成电路创新创业大赛\par}
  \vspace{2.6cm}
  {\zihao{2}\bfseries 基于 RISC-V 指令集的 CPU 设计\par}
  \vspace{0.8cm}
  {\zihao{2}\bfseries 后续优化说明报告\par}
  \vfill
  \begin{tabular}{rl}
    \zihao{4}\bfseries 参赛杯赛： & \zihao{4}“竞业达”企业命题 \\
    \zihao{4}\bfseries 队伍编号： & \zihao{4}CICC1004252 \\
    \zihao{4}\bfseries 团队名称： & \zihao{4}和溢位 \\
  \end{tabular}
  \vspace{2cm}
\end{titlepage}

\tableofcontents
\clearpage

\section{报告说明}

分赛区决赛主报告完成后，我们将主要精力转向处理器性能与 FPGA 时序优化。此时数字孪生平台已经关闭，
后续工作主要通过 Verilator、Vivado XSim 和实现时序报告验证；
由于数字孪生平台的关闭，缺乏上板验证的依据，而主报告内容则均经过正式上板验证，我们选择将后续
优化单独写成本报告，作为对主报告的补充。

主报告所描述的设计基于较旧的版本，但架构没有改变。
本报告将不再重复主报告中的完整架构等信息，主要补充说明此后完成的主要优化及量化结果。

优化始终保持\textbf{顺序单发射}的总体架构。我们通过性能计数器、周期级 RTL 仿真、Vivado
布局布线结果指导优化：先定位主要停顿来源，再用小范围改动验证假设，最后同时检查功能、周期以及
建立/保持时序。

性能数据来自 Verilator 周期级 RTL 仿真，并按对应实现频率换算为评估时间。
仅有仿真收益而无法满足时序的方案不作为正式结果。

本报告保留两套性能数据：\code{withMext-v2} 用于贯穿主报告版本至最终版本的统一 RTL 仿真复测；
西北赛区样例用于呈现后期优化对更大程序及不同访存、分支特征的效果。前者每次完整运行均提交
380,344,384 条动态指令，表中时间由完整周期数除以 280 MHz 得到；后者约含 14.47 亿条动态指令。
两套程序的运行时间\textbf{不能横向比较}，本报告只比较同一程序内部相邻版本的变化。

\section{整体优化思路}

\begin{enumerate}
  \item \textbf{用数据定位瓶颈。} 软件仿真快速统计、添加更多性能计数器。
  \item \textbf{优先减少高频停顿。} 主要针对 \code{load-use} 相关和乘法器。
  \item \textbf{控制组合路径规模。} 缩小旁路范围、减少宽 Mux、局部化地址生成和简化接口。
  \item \textbf{固定项目实现。} 对 RTL、IP 参数配置和 Vivado 项目统一纳入 Git 仓库管理。
\end{enumerate}

\section{主要优化内容}

\subsection{早期性能优化}

主报告版本后，我们先围绕乘法、方向预测、缓存写入和 load-use 停顿开展优化：

\begin{itemize}
  \item \textbf{缩短乘法流水延迟。} 将乘法器从 6 周期逐步缩短并验证 5、4、3 周期方案。4 周期方案
  在性能与时序之间取得较好平衡；3 周期虽然更快，但产生大范围时序违例，因而没有作为最终方案。
  \par\noindent\textbf{支撑：}{\kaishu 性能计数器在完整程序中记录到约 1050 万次普通乘法，说明每减少一拍都可能形成稳定收益。}
  \item \textbf{改进分支方向预测。} 为条件分支加入 2 bit 饱和计数器。该程序原有预测准确率已经较高，
  因此收益不大，但改动能够保持功能和时序正确。
  \par\noindent\textbf{支撑：}{\kaishu withMext-v2 完整仿真记录的预测准确率约为 98.75\%，与该项收益较小的现象一致；
  }
  \item \textbf{优化 DCache 的整字写行为。} 原设计在 store 后使缓存项失效，导致热点中的后续 load
  反复未命中。新方案对可缓存的整字 \code{sw} 执行 write-allocate，同时用写入优先级和 epoch 防止
  较早的 load fill 覆盖新数据；byte/half store 仍采取保守失效策略。
  \par\noindent\textbf{支撑：}{\kaishu 计数器记录到约 2145 万次可缓存整字 store，且相邻版本时间从
  1.943 s 降至 1.868 s。}
  \item \textbf{优化 load-use 数据相关。} 性能计数表明主要瓶颈不是分支，而是 load 后紧邻的
  \code{ADD/ADDI}。晚数据旁路允许这类消费者提前进入 EXU，在缓存命中数据到达时直接替换源操作数；
  若数据尚未就绪则保持指令，避免使用无效数据。
  \par\noindent\textbf{支撑：}{\kaishu 完整计数器分别观察到 57,668,883 次 rs1 和 5,312,000 次 rs2
  的相邻 load--ADD/ADDI 机会；对应仿真时间从 1.868 s 降至 1.643 s。}
  \par\noindent\textbf{说明：}我们把适用范围限制为容易验证的加法消费者，并将普通加法结果与重定向
  路径拆开，避免 DCache 命中信号经 ready 链一路影响 IFU，从而控制新增旁路的时序代价。
  \item \textbf{提高工程可复现性。} 修复 Vivado 工程曾读取旧目录生成 RTL 的路径问题，统一项目内
  打包源、IP 输出物、时钟参数和实现策略，确保报告的物理结果确实对应当前候选。
\end{itemize}

\subsection{基于西北赛区样例的后续优化}

这一阶段的优化方向来自西北赛区样例的完整性能统计、NEMU 模型和固定 RTL 前缀。两套样例的汇总结果
在后文分别展示。

\begin{itemize}
  \item \textbf{练习 B 扩展功能。} 尝试练习实现了一些 RISC-V B 扩展功能。较复杂的操作采用多周期
  执行与明确的 ready/valid 控制，在保证功能的同时限制单周期逻辑深度。
  \item \textbf{重新收敛 280 MHz 时序。} load-use 候选在 280 MHz 下仍有 WNS $-0.084$ ns 违例。
  我们把跨模块的宽控制改为粗粒度类别，在多周期单元内部完成具体译码；同时为 LSU 地址生成设置独立
  的 18 bit 通路，把旁路选择移到地址加法进位链之后，缩短“选择—加法—地址输出”的串联路径。
  \par\noindent\textbf{支撑：}{\kaishu 对应实现把 WNS 从 $-0.084$ ns 恢复到 $0.000$ ns；同一完整
  仿真的周期数保持不变，说明这轮调整解决的是组合路径，而不是以增加停顿换取时序。}
  \par\noindent\textbf{说明：}在 RTL 结构缩短后，再用 post-route 物理优化完成最后的布线整理。
  \item \textbf{改进乘除法实现。} 低位乘法使用更短流水，乘法结果选择从通用结果 Mux 中局部化；除法
  采用 32 轮恢复除法，每拍完成一次移位、比较和减法，并用明确的 ready/valid 握手控制流水停顿。
  最后一轮的符号修正后来单独分拍，避免 33 bit 减法与 32 bit 取负串在同一关键路径上。
  \par\noindent\textbf{支撑：}{\kaishu M 扩展定向测试全部通过；符号修正分拍在固定前缀中只增加
  12 周期，并移除了原除法结束路径对应的时序违例族。}
  \item \textbf{提升窄加载和缓存局部性。} 这里的 DCache shadow 是主缓存状态的一份轻量快速副本，
  保存快速命中判断所需的 valid、tag 和数据字节；它不是新的一级或二级缓存。设计保证 shadow 与主
  缓存一致：load fill 时同步更新，store 命中时按写掩码更新，无法安全合并时保守失效。
  \par\noindent\textbf{支撑：}{\kaishu 西北赛区样例的完整模型找到 87,172,384 个可节省周期，并指出
  热点集中在约 7.36 亿条指令之后，短前缀不足以代替完整运行。}
  \par\noindent\textbf{说明：}在 shadow 上直接完成 \code{LBU/LHU} 的字节、半字提取，可让紧邻的
  简单消费者更早获得正确数据，同时避免旧数据被错误旁路。
  \item \textbf{保留部分写命中的缓存行。} 对能够安全合并的部分 store，不再像旧设计那样使整条
  shadow 缓存行失效，而是按照 byte-enable 只更新被写入的字节并保留其余字节。这样既维持了 shadow
  与主存写入的一致性，又减少了随后读取同一缓存行时重新 miss 的概率。
  \par\noindent\textbf{支撑：}{\kaishu 完整缓存模型的总命中率由 86.0341\% 提高到 95.8535\%，预测
  节省 298.394 ms；实际记录由 6.481 s 降至 6.182 s，差值为 299 ms。}
  \item \textbf{扩大并优化 BTB。} 将 BTB 的有效物理项扩展为 32 项，并把更新 shadow 放入 LUTRAM，
  减少容量冲突。主查询存储保留预测目标等完整表项；较窄的更新 shadow 仅保存 valid、tag、分支类型
  和 2 bit 方向计数器，使更新端无需绕回主查询数据通路。将它改为 LUTRAM 后，预测语义不变，但减少
  触发器、时钟扇出和布线压力。
  \par\noindent\textbf{支撑：}{\kaishu 模型预计 32 项 BTB 少 4,187,905 次 miss；RTL A/B 标定每次
  miss 为 3 周期，换算 44.870 ms，与 6.553 s 到 6.508 s 的 45 ms 变化一致。LUTRAM 化前后固定前缀
  周期完全相同。}
  \item \textbf{识别并预测标准返回。} RISC-V 的标准 \code{ret} 是伪指令，其精确机器码对应
  \code{jalr x0, 0(x1)}，即 \code{0x00008067}。译码级直接比较这一完整指令编码：只有完全匹配时，
  才把 BTB 中已记录的目标作为无条件预测目标；其他 JALR 仍由原执行级路径解析。这样不需要增加通用
  JALR 的 32 bit 目标比较器，也不会改变非标准间接跳转的执行语义。
  \par\noindent\textbf{支撑：}{\kaishu 模型预计减少 2,575,420 次 miss，即 27.594 ms；记录由
  6.508 s 降至 6.481 s。}
  \item \textbf{缩短关键接口并分拍除法修正。} 在部分写命中版本之后，我们继续缩短返回地址
  旁路：当 \code{ret} 依赖 EXU 正在产生的 \code{x1} 时，不再要求同拍完成预测目标比较，而是保守地
  交给原重定向路径处理；同时删除重复地址字段，将跨模块数据缩窄到实际使用的 22 bit。随后把同拍
  ALU 旁路从缓存/LSU/WBU 的通用选择树中分离，把 WBU 到地址生成器的专用旁路缩减为 22 bit 数据和
  一个有效位，并把 DCache 查询接口从完整 32 bit 地址缩减为实际需要的 index 和 tag。这些接口调整
  在本次完整仿真中均未增加周期。

  这一项中唯一有明确周期代价的是迭代除法器分拍。原先最后一轮迭代需要在同一拍完成 33 bit 减法和
  32 bit 符号修正；新设计增加 \code{correct} 状态，先寄存无符号商和余数，下一拍再做可选的二进制
  补码修正。因此，每次普通除法或求余多占用一拍。本次完整仿真中共有 12 次此类操作，因此总周期
  精确增加 12，即在 280 MHz 下约 $0.043\,\mu$s。
  其余接口缩减均未改变周期数。上述 RTL 修改将违例缩小到约 $-0.004$ ns。
  \item \textbf{选择匹配网表的实现策略。} 在完成上述 RTL 结构调整后，再使用 Auto/Explore、
  NoTimingRelaxation 和 post-route 物理优化关闭最后几皮秒的布线违例，最终得到 WNS $+0.013$ ns。
  因此最终无违例结果来自“RTL 关键路径缩减、除法修正分拍、实现策略选择”的组合，而不是单独
  更换实现策略；这也解释了时序重新收敛时出现的 12 周期代价。
\end{itemize}

\subsection{性能结果}

\begin{table}[H]
  \centering
  \caption{统一采用 withMext-v2 的完整 RTL 仿真性能演进}
  \renewcommand{\arraystretch}{1.45}
  \setlength{\tabcolsep}{5pt}
  \resizebox{\textwidth}{!}{%
  \begin{tabular}{@{}lcccc@{}}
      \toprule
      优化里程碑 & 频率 & 仿真时间 & 较前基线提升 & 时序状态 \\
      \midrule
      主报告版本 & 280 MHz & 2.020 s & --- & WNS $+0.050$ ns \\
      乘法与方向预测优化 & 280 MHz & 1.943 s & $3.8\%$ & WNS $+0.031$ ns \\
      DCache 整字写分配 & 280 MHz & 1.868 s & $3.9\%$ & WNS $+0.019$ ns \\
      load-use 晚数据旁路 & 280 MHz & 1.643 s & $12.0\%$ & WNS $-0.084$ ns \\
      迭代 RTL 除法器 & 280 MHz & 1.606910 s & $2.2\%$ & 通过 \\
      LSU 地址通路与 280 MHz 收敛 & 280 MHz & 1.606910 s & $0.0\%$ & WNS $0.000$ ns \\
      窄加载与 DCache shadow & 280 MHz & 1.607404 s & $-0.031\%$ & WNS $0.000$ ns \\
      32 项 BTB 与更新 shadow & 280 MHz & 1.607403 s & $+0.000055\%$ & WNS $-0.011$ ns \\
      精确 ret 识别与预测 & 280 MHz & 1.607403 s & $+0.000013\%$ & WNS $-0.091$ ns \\
      部分写命中保留 & 280 MHz & 1.607403 s & $+0.000001\%$ & WNS $-0.072$ ns \\
      接口缩减与除法修正分拍 & 280 MHz & 1.607403 s & $-0.000003\%$ & WNS $-0.004$ ns \\
      \textbf{最终时序闭合版本} & \textbf{280 MHz} & \result{1.607403 s} & $0.0\%$ & WNS $+0.013$ ns \\
      \bottomrule
  \end{tabular}}
\end{table}

最终版本完整运行 450,072,868 周期、380,344,384 条指令，CPI 为 1.183330。相对主报告版本的
2.020 s，统一口径下仿真时间缩短约 \result{20.4\%}。后期若干优化主要用于减少资源与关闭时序，
在该自测程序中的周期变化很小；表中仍保留精确百分比，避免把时序优化误写成性能收益。最终实现的
建立时间 WNS 为 $+0.013$ ns、保持时间 WHS 为 $+0.072$ ns。

\subsection{西北赛区样例结果}

西北赛区样例与 \code{withMext-v2} 的动态指令组成不同，以下数据仅在本表内部逐行比较。表中里程碑
与前述优化项一一对应；LUTRAM shadow、返回旁路和接口缩减主要改善时序，周期不变或只有已量化的小代价。

\begin{table}[H]
  \centering
  \caption{西北赛区样例的性能与时序演进}
  \setlength{\tabcolsep}{5pt}
  \begin{tabular}{@{}lccc@{}}
    \toprule
    优化里程碑 & 评估时间 & 较前基线提升 & 时序状态 \\
    \midrule
    多周期单元与 280 MHz 收敛 & 6.866 s & --- & WNS $0.000$ ns \\
    窄加载与 DCache shadow & 6.553 s & $4.6\%$ & WNS $0.000$ ns \\
    32 项 BTB 与更新 shadow & 6.508 s & $0.7\%$ & WNS $-0.011$ ns \\
    精确 ret 识别与预测 & 6.481 s & $0.4\%$ & WNS $-0.091$ ns \\
    部分写命中保留 & 6.182 s & $4.6\%$ & WNS $-0.072$ ns \\
    \textbf{接口整理与最终时序闭合} & \result{6.184 s} & $-0.03\%$ & WNS $+0.013$ ns \\
    \bottomrule
  \end{tabular}
\end{table}

最终结果相对 6.866 s 基线缩短约 \result{9.9\%}。6.182 s 到 6.184 s 的小幅回退出现在除法修正
分拍和接口整理之后；固定前缀明确观察到 12 次除法各增加一拍，而最终版本同时将 WNS 从负值收敛到
$+0.013$ ns，因此选择 6.184 s 的时序闭合版本作为最终结果。

\subsection{性能与 WNS 走向}

图中蓝线为同一样例内部的评估时间，红色虚线为对应实现的 WNS；灰色零线以上表示建立时序通过。
纵轴分别缩放，折线用于观察趋势，不用于比较两套样例的绝对时间。

\begin{figure}[H]
  \centering
  \begin{tikzpicture}
    \begin{axis}[
      width=0.835\textwidth,height=6.2cm,
      symbolic x coords={主报告,乘分支,DCache,晚旁路,LSU地址通路,Shadow,BTB,Ret,部分写,最终},
      xtick=data,x tick label style={rotate=32,anchor=east,font=\scriptsize},
      ylabel={评估时间（s）},ymin=1.55,ymax=2.05,
      axis y line*=left,axis x line*=bottom,
      grid=major,legend style={at={(0,1.02)},anchor=south west,font=\scriptsize}]
      \addplot[blue!75!black,mark=*] coordinates {
        (主报告,2.020) (乘分支,1.943) (DCache,1.868) (晚旁路,1.643)
        (LSU地址通路,1.606910) (Shadow,1.607404) (BTB,1.607403) (Ret,1.607403)
        (部分写,1.607403) (最终,1.607403)};
      \addlegendentry{性能}
    \end{axis}
    \begin{axis}[
      width=0.835\textwidth,height=6.2cm,
      symbolic x coords={主报告,乘分支,DCache,晚旁路,LSU地址通路,Shadow,BTB,Ret,部分写,最终},
      xtick=\empty,ylabel={WNS（ns）},ymin=-0.11,ymax=0.07,
      axis y line*=right,axis x line=none,
      legend style={at={(1,1.02)},anchor=south east,font=\scriptsize}]
      \draw[gray,dashed] (rel axis cs:0,0.6111) -- (rel axis cs:1,0.6111);
      \addplot[red!75!black,dashed,mark=square*] coordinates {
        (主报告,0.050) (乘分支,0.031) (DCache,0.019) (晚旁路,-0.084)
        (LSU地址通路,0.000) (Shadow,0.000) (BTB,-0.011) (Ret,-0.091)
        (部分写,-0.072) (最终,0.013)};
      \addlegendentry{WNS}
    \end{axis}
  \end{tikzpicture}
  \caption{withMext-v2 的性能与 WNS 走向}
\end{figure}

\begin{figure}[H]
  \centering
  \begin{tikzpicture}
    \begin{axis}[
      width=0.835\textwidth,height=6.0cm,
      symbolic x coords={时序基线,窄加载,BTB,Ret,部分写,最终},
      xtick=data,x tick label style={rotate=25,anchor=east,font=\scriptsize},
      ylabel={评估时间（s）},ymin=6.05,ymax=6.95,
      axis y line*=left,axis x line*=bottom,
      grid=major,legend style={at={(0,1.02)},anchor=south west,font=\scriptsize}]
      \addplot[blue!75!black,mark=*] coordinates {
        (时序基线,6.866) (窄加载,6.553) (BTB,6.508) (Ret,6.481) (部分写,6.182) (最终,6.184)};
      \addlegendentry{性能}
    \end{axis}
    \begin{axis}[
      width=0.835\textwidth,height=6.0cm,
      symbolic x coords={时序基线,窄加载,BTB,Ret,部分写,最终},
      xtick=\empty,ylabel={WNS（ns）},ymin=-0.11,ymax=0.04,
      axis y line*=right,axis x line=none,
      legend style={at={(1,1.02)},anchor=south east,font=\scriptsize}]
      \draw[gray,dashed] (rel axis cs:0,0.7333) -- (rel axis cs:1,0.7333);
      \addplot[red!75!black,dashed,mark=square*] coordinates {
        (时序基线,0.000) (窄加载,0.000) (BTB,-0.011) (Ret,-0.091)
        (部分写,-0.072) (最终,0.013)};
      \addlegendentry{WNS}
    \end{axis}
  \end{tikzpicture}
  \caption{西北赛区样例的性能与 WNS 走向}
\end{figure}

\section{验证与取舍}

优化过程中同时保留成功和失败实验，避免只依据最好的一次数据选择方案。功能验证覆盖基础 CPU 测试、
load-store、分支与递归测试、M 扩展定向架构测试，以及相关功能练习测试；涉及 CSR 的候选还运行
RT-Thread 至命令行提示符。性能候选使用完整程序或固定代表性前缀核对周期趋势，并通过 Vivado 实现
结果检查目标频率下的建立与保持时序。

若某方案明显减少周期但产生大范围负 WNS，例如把 DCache 查询前移到 IDU、使用宽通用队列、过度扩展
组合旁路或采用 3 周期乘法器，我们会保留实验结论但回退实现。轻微负 WNS 的候选用于分析优化方向和
呈现演进过程，不作为最终闭合版本。最终版本还要求 PLL 请求频率与实际生成频率一致，且 setup/hold 均通过。

\section{总结}

在不改变顺序单发射总体架构的前提下，后续优化把工作重点从“完成功能”推进到“依据真实热点降低周期，
并让优化后的网表稳定通过 FPGA 时序”。统一采用 \code{withMext-v2} 完整仿真后，280 MHz 换算时间
由主报告版本的 2.020 s 降低到最终版本的 1.607403 s，缩短约 20.4\%；西北赛区样例在自身统一口径下
由 6.866 s 降低到最终时序闭合版本的 6.184 s，缩短约 9.9\%。两组数据相互独立，但共同说明优化方向
在不同程序特征下具有可观察效果。

这些工作不仅包含乘法、缓存、旁路、分支预测和相关功能练习，也包括 IP 接口核对、RTL 除法器、
Vivado 输入固定和实现策略选择。最终结果表明，处理器在功能正确的基础上进一步提高了有效吞吐率，
并保持了可复现的实现时序。

\appendix
\section{关键优化提交索引}

下表给出正文关键里程碑对应的完整 Git commit hash。部分里程碑由多个小提交共同构成；表中按功能边界
列出，便于结合提交说明和 RTL diff 复核。

\small
\begin{longtable}{@{}p{0.43\textwidth}p{0.52\textwidth}@{}}
  \caption{关键优化与 Git 提交对应关系}\\
  \toprule
  关键内容 & 完整 commit hash \\
  \midrule
  \endfirsthead
  \toprule
  关键内容 & 完整 commit hash \\
  \midrule
  \endhead
  主报告性能基线 & \ttfamily d41aa7495b7b5a19f666ad627f998debb232b5b0 \\
  乘法、方向预测与整字写分配组合 & \ttfamily 196077bdd8187c62e6e9988820ce27fffc87247e \\
  load-use 晚数据旁路基线 & \ttfamily d8364daf1df0ea203ab661bf5c093e208b9414f9 \\
  迭代 RTL 除法器 & \ttfamily 85d3dad160c1dfcc7ef200fc14da22b76319dc77 \\
  LSU 地址旁路后移 & \ttfamily 11f93f5f4cb41e637056fc7b341461bc13a11af2 \\
  缓存窄加载 shadow & \ttfamily e71bb475819cfcec19bde6553f54ab37091db591 \\
  窄 store 的 shadow 一致性修复 & \ttfamily 5a73acc1b28c991076a1804d434c8a9e4448ad63 \\
  热点 LBU/LHU 提取通路 & \ttfamily b5731d80617cf29ff2d8a76d71f8c64c6dca0443 \\
  多周期单元内部精确译码 & \ttfamily d71fa1a364500fe00c27440e821ee80d92d0b743 \\
  有符号除法溢出迭代处理 & \ttfamily 906249b9443b2b25ee1e064f9f3ab296d75d9eea \\
  启用 32 项 BTB & \ttfamily 8eb6b8ccd5538f099765e664b0bab63013c483a8 \\
  BTB 更新 shadow 改用 LUTRAM & \ttfamily a72b1f9c2756ccf211eb8bb897b06b1acb3dbf7f \\
  仅对精确 ret 使用间接预测 & \ttfamily f5143a0c398e9f4cb0f2b3c1e7d34b286b5a179b \\
  保留部分写命中的缓存行 & \ttfamily d205c9aa7a41c8d8611033fd6aee0573fb863014 \\
  缩短返回地址旁路 & \ttfamily 3fd8030b5e0ef0028e31df3fcf3dd0737f090eba \\
  缩短 ALU 旁路并分拍除法修正 & \ttfamily c4279e438726cceae21fc71321ee5ee306bb95e9 \\
  缩窄 WBU 地址旁路 & \ttfamily e3a7baaffdf632a1bddd790ae8cedb7df9fcf9cd \\
  缩窄 DCache 地址接口 & \ttfamily 9343d53069f6bb268dd1210ae2b083a70699008f \\
  最终冻结 RTL & \ttfamily 58f371d5dbff5e6a8c730ab59ea6d39048012159 \\
  最终 Vivado 输入恢复 & \ttfamily c6e15a9d0e43099bea863c45e5f0020a5bac597a \\
  \bottomrule
\end{longtable}
\normalsize

\end{document}
